\section{Introducción}
\pendiente{Presentar el problema, importancia de la cinemática de un robot antropomórfico y propósito de la simulación. Sustentar el contexto con citas.}
\textbf{Objetivo general propuesto:} modelar y simular en Python la cinemática de un manipulador antropomórfico de 3 GDL basado en un robot documentado, integrando análisis directo, inverso, diferencial, singularidades y control cinemático.
\textbf{Objetivos específicos:} establecer DH y cinemática directa; obtener soluciones inversas; derivar y verificar Jacobianos; caracterizar singularidades de posición; evaluar una ley de control mediante métricas reproducibles.
\textbf{Alcance provisional:} geometría y cinemática rígida nominal del robot de \textcite{ramzy2024}. Confirmar si satisface la consigna \parencite{consigna2026}; no afirmar simulación de flexibilidad ni réplica experimental.
